\section{Resultados}

A análise compara as 53 audiências relacionadas a grupos minoritários
(grupo M) às 153 demais audiências do corpus (grupo C). Em conjunto, os
resultados indicam diferenças mais consistentes no conteúdo das justificações
e nas manifestações explícitas de respeito do que nas condições formais de
participação, no nível médio de justificação ou na cobertura institucional. A
interpretação dos achados deve permanecer exploratória, pois parte das
medidas deriva de classificações por modelos de linguagem ainda sem validação
humana sistemática, e a própria rotulagem do recorte M/C está em processo de
conferência.

\subsection{Interrupções e condições de participação}

A análise de interrupções considerou 4.570 falas elegíveis, classificadas
segundo o papel do falante na audiência. Nas audiências do grupo M, 65 de 751
falas de convidados foram marcadas como interrompidas, em comparação com 5 de
113 falas de parlamentares. No grupo C, a mesma proporção correspondeu a 314
de 2.715 falas de convidados e 87 de 991 falas de parlamentares.

A regressão logística com GEE, agrupando as falas por audiência, não encontrou
evidência estatisticamente significativa de que convidados fossem interrompidos
com frequência distinta da observada para parlamentares nas audiências do grupo
C ($OR=1{,}45$, IC95\% $[0{,}81;2{,}59]$, $p=0{,}211$). O termo de interação
entre participante convidado e audiência do grupo M apresentou $OR=1{,}96$
(IC95\% $[0{,}80;4{,}79]$, $p=0{,}139$). Embora a estimativa pontual esteja na
direção de uma diferença mais acentuada nas audiências sobre minorias, o
intervalo de confiança é amplo e inclui a ausência de efeito.

Dessa forma, os resultados atuais não sustentam a hipótese de que convidados
sejam desproporcionalmente mais interrompidos em audiências relacionadas a
grupos minoritários. Esse achado deve ser interpretado com cautela, uma vez
que a regra utilizada para identificar interrupções ainda se encontra em
revisão.

\subsection{Conteúdo das justificações}

A diferença mais clara entre os grupos aparece no conteúdo das justificações.
Nas audiências do grupo M, a mediana da proporção de justificações orientadas
ao bem comum sensível à diferença foi de 18,18\%, frente a 5,43\% no grupo C.
A diferença entre as medianas foi de 12,75 pontos percentuais (IC95\%
$[7{,}78;17{,}43]$), com $p<0{,}001$ e efeito rank-biserial de $r=0{,}63$.

Esse resultado sugere que, nas audiências relacionadas a grupos minoritários,
as justificações classificadas pelo modelo recorrem com maior frequência a
formulações do bem comum que reconhecem explicitamente diferenças entre grupos,
desigualdades ou condições sociais específicas.

Em contraste, não foi observada diferença no uso de justificações orientadas ao
interesse de grupo. A mediana foi igual a zero nos dois conjuntos de
audiências, com $p=0{,}487$ e $r=-0{,}06$. Assim, os dados sustentam apenas a
primeira parte da hipótese relativa ao conteúdo das justificações: há maior
presença de bem comum sensível à diferença nas audiências M, mas não menor
incidência de justificações de interesse de grupo.

\subsection{Nível de justificação}

O nível médio de justificação apresentou distribuições semelhantes nos dois
grupos. Em uma escala de 0 a 3, a mediana foi de 0,96 nas audiências M e 1,00
nas audiências C. A diferença estimada entre as medianas foi de $-0{,}04$
(IC95\% $[-0{,}13;0{,}01]$), com $p=0{,}115$ e $r=-0{,}15$.

Portanto, os resultados não fornecem evidência de uma diferença sistemática no
nível de elaboração das justificações entre audiências relacionadas a minorias
e as demais audiências do corpus.

\subsection{Respeito e hostilidade}

Foram observadas diferenças consistentes nas manifestações explícitas de
respeito. A mediana da proporção de códigos de respeito classificados como
explícitos positivos foi de 26,67\% nas audiências do grupo M e 19,44\% no
grupo C. A diferença entre as medianas foi de 7,22 pontos percentuais (IC95\%
$[5{,}49;9{,}52]$), com $p<0{,}001$ e $r=0{,}59$.

O nível médio de respeito também foi superior no grupo M: a mediana foi de
0,579 nas audiências M e 0,543 nas audiências C, em uma escala de 0 a 1. A
diferença entre as medianas foi de 0,036 (IC95\% $[0{,}025;0{,}046]$), com
$p<0{,}001$ e $r=0{,}56$.

Esse padrão, entretanto, não foi acompanhado por maior hostilidade. Ocorrências
de respeito negativo ou degradante foram identificadas em 4 das 53 audiências
do grupo M e em 22 das 153 audiências do grupo C. O teste exato de Fisher não
indicou diferença estatisticamente significativa ($OR=0{,}49$, $p=0{,}237$).

Assim, a hipótese originalmente formulada, segundo a qual audiências sobre
minorias combinariam maior elogio e maior hostilidade sem alteração do nível
médio de respeito, não é sustentada integralmente. Os dados indicam maior
manifestação explícita de respeito nas audiências M, mas não maior ocorrência
de hostilidade.

\subsection{Técnicas de persuasão}

A análise de persuasão foi realizada em uma subamostra balanceada de 20
audiências: 10 relacionadas a minorias e 10 de temáticas variadas. A
classificação considerou seis categorias de técnicas persuasivas --- ataque à
reputação, justificativa, simplificação, distração, chamada à ação e linguagem
manipulativa --- além da categoria ``nenhuma'', que indica a ausência dessas
técnicas no trecho classificado.

Antes da correção para múltiplas comparações, foram observadas diferenças
entre os grupos em ataque à reputação, justificativa, chamada à ação, linguagem
manipulativa e ausência de técnica. Após a aplicação da correção de Holm para
as sete categorias, apenas a categoria ``nenhuma'' permaneceu abaixo do limiar
de 0,05. A proporção média de trechos sem nenhuma técnica foi 7,93 pontos
percentuais menor no grupo M ($p_{\mathrm{Holm}}=0{,}047$).

As demais diferenças apontaram para maior uso de recursos persuasivos nas
audiências M, mas não permaneceram estatisticamente significativas após a
correção. A chamada à ação apresentou diferença de 9,58 pontos percentuais
($p_{\mathrm{Holm}}=0{,}072$); justificativa, de 7,43 pontos percentuais
($p_{\mathrm{Holm}}=0{,}083$); linguagem manipulativa, de 6,15 pontos
percentuais ($p_{\mathrm{Holm}}=0{,}083$); e ataque à reputação, de 8,30
pontos percentuais ($p_{\mathrm{Holm}}=0{,}116$).

Esses resultados sugerem, de forma preliminar, maior densidade de recursos
persuasivos nas audiências do grupo M. Contudo, a evidência corrigida não
permite atribuir essa diferença de modo robusto a uma técnica específica. A
pequena subamostra utilizada nessa análise e a necessidade de validação humana
das classificações recomendam cautela adicional na interpretação.

\subsection{Participação discursiva e cobertura institucional}

A comparação entre participação discursiva e presença na matéria institucional
não revelou diferença significativa entre os grupos. O indicador de déficit da
sociedade civil mede a diferença entre a proporção de palavras ditas por
convidados e a proporção de posições atribuídas a convidados na matéria da
Agência Câmara. Valores positivos indicam que a sociedade civil falou mais do
que apareceu na cobertura.

A mediana do indicador foi de 1,20 ponto percentual nas audiências M e de 4,19
pontos percentuais nas audiências C. A diferença estimada entre as medianas
(M -- C) foi de $-3{,}00$ pontos percentuais (IC95\% $[-9{,}89;4{,}56]$), com
$p=0{,}470$ e $r=-0{,}07$.

Como audiências maiores podem produzir padrões distintos de seleção na matéria
jornalística, estimou-se adicionalmente um modelo controlando o tamanho da
sessão. O efeito associado ao grupo M permaneceu pequeno e não significativo
($\beta=-0{,}021$, $p=0{,}610$). Portanto, os resultados atuais não sustentam a
hipótese de que a discrepância entre espaço discursivo e representação
institucional seja sistematicamente maior nas audiências relacionadas a
minorias.

No corpus como um todo, entretanto, a comparação nominal entre participantes
da audiência e pessoas mencionadas na matéria revela uma assimetria substantiva
de cobertura: 1.732 de 2.543 falantes identificados, aproximadamente 68,1\%,
não aparecem na matéria institucional associada. Esse resultado caracteriza uma
propriedade geral do processo de síntese institucional, embora não diferencie,
por si só, audiências M e C.
